\section{Conclusion}
\label{sec:conclusion}

We introduced Adaptive Dual-Branch Decoding, a training-free method for reducing hallucinations in LVLMs. ADBD replaces ONLY's frozen one-layer binary decision with a continuous TVER state refined across selected decoder layers by head-specific adaptive gains. The state produces complementary high- and low-TVER auxiliary branches, which are independently reinforced or penalized according to their disagreement with the normal distribution. An anchored three-branch rule combines these signals without requiring a separate full LVLM invocation.

On POPE with LLaVA-1.5-7B, ADBD achieves the strongest performance among the evaluated decoding configurations. \TODO{Insert exact macro and split-specific improvements.} On MME Hallucination, the method improves Position and Color but decreases Existence and Count, yielding a reported aggregate score of 620.00 compared with 635.55 for ONLY; the aggregate must be verified because the supplied ADBD category scores sum to 630.00. Ablations attribute the POPE improvement to \TODO{verified component findings}, while branch-level analysis shows \TODO{verified decoding-dynamics finding}. Overall, the results suggest that complementary and adaptively accumulated attention evidence can provide a more informative contrastive signal than a static single-layer mask, but its gains are category-dependent rather than universal.
